\apartado{2.9}{Distribución normal}{distrib}

\bajada{La curva que permite evaluar cualquier punto de corte sin volver a contar el archivo. Cierra la unidad respondiendo dónde conviene ponerlo.}
\ejercicio{0}{%
  Según la regla 68–95–99,7, ¿qué porcentaje de los casos cae entre
  $\mu-\sigma$ y $\mu+\sigma$? ¿Y fuera de $\mu\pm 2\sigma$?}{}
\renglones[3]
\respuesta{68\% dentro de $\mu\pm\sigma$; 5\% fuera de $\mu\pm 2\sigma$.}
\solucion{%
  Entre $\mu\pm\sigma$ cae el \textbf{68\%}.
  \\[0.3em]
  Fuera de $\mu\pm 2\sigma$ cae el $100-95 = \mathbf{5\%}$, repartido en las
  dos colas: un 2,5\% de cada lado.}

\ejercicio{1}{%
  Un test de inteligencia tiene $\mu=100$ y $\sigma=15$.
  ¿Entre qué dos valores cae el 68\% de la población? ¿Y el 95\%?}{}
\renglones[3]
\respuesta{68\%: entre 85 y 115. 95\%: entre 70 y 130.}
\solucion{%
  68\%: $100 \pm 15$, o sea entre \textbf{85 y 115}.
  \\[0.3em]
  95\%: $100 \pm 2(15) = 100\pm 30$, o sea entre \textbf{70 y 130}.}

\ejercicio{1}{%
  Con el mismo test ($\mu=100$, $\sigma=15$), calcula la puntuación $z$ de
  alguien que obtuvo 130. Interpreta el resultado en una frase.}{}
\renglones[3]
\respuesta{$z=2$: está a 2 desviaciones estándar por encima del promedio.}
\solucion{%
  \[ z = \frac{130-100}{15} = \frac{30}{15} = 2 \]
  Está a \textbf{2 desviaciones estándar por encima} del promedio. Por la
  regla 68–95–99,7, sólo un 2,5\% de la población supera ese valor.}

\ejercicio{2}{%
  Con el mismo test, ¿qué puntaje bruto corresponde a $z=-1{,}5$?
  ¿Y qué $z$ le corresponde a un puntaje de 92?}{}
\renglones[3]
\respuesta{$z=-1{,}5 \rightarrow 77{,}5$; \; $x=92 \rightarrow z\approx -0{,}53$.}
\solucion{%
  Despejando $x$ de $z=\dfrac{x-\mu}{\sigma}$:
  \[ x = \mu + z\sigma = 100 + (-1{,}5)(15) = 100 - 22{,}5 = 77{,}5 \]
  \\[0.3em]
  Y al revés:
  \[ z = \frac{92-100}{15} = \frac{-8}{15} \approx -0{,}53 \]
  \\[0.3em]
  El signo dice de qué lado del promedio está; el valor absoluto, cuán lejos.}

\ejercicio{2}{%
  En una distribución normal con $\mu=50$ y $\sigma=10$, ¿qué proporción de
  los casos supera el valor 70? Resuélvelo con la regla 68–95–99,7, sin tabla.}{}
\renglones[3]
\respuesta{$z=2$, así que aproximadamente el 2,5\%.}
\solucion{%
  \[ z = \frac{70-50}{10} = 2 \]
  Fuera de $\mu\pm 2\sigma$ queda el 5\%, repartido en dos colas simétricas.
  La cola de arriba se lleva la mitad: \textbf{2,5\%}.}

\ejercicio{3}{%
  En las 200 fichas, el puntaje de depresión tiene $\mu=6{,}32$ y
  $\sigma=4{,}64$. El corte del tamizaje está en 10.
  \begin{enumerate}[label=(\alph*),topsep=2pt,itemsep=1pt]
    \item Calcula la puntuación $z$ del corte.
    \item Usando la tabla normal, $P(Z \ge 0{,}79) \approx 0{,}214$.
          ¿Qué significa ese número?
    \item Compara con la proporción real de positivos del archivo ($43/200$).
  \end{enumerate}}{}
\renglones[5]
\respuesta{(a) $z\approx 0{,}79$. (b) 21,4\% superaría el corte según el
modelo. (c) Real: 21,5\%. Casi idéntico.}
\solucion{%
  (a) \[ z = \frac{10-6{,}32}{4{,}64} = \frac{3{,}68}{4{,}64} \approx 0{,}79 \]
  \\[0.3em]
  (b) Significa que, si los puntajes siguieran una distribución normal, el
  \textbf{21,4\%} de las personas superaría el corte de 10.
  \\[0.3em]
  (c) La proporción real es $\dfrac{43}{200} = 21{,}5\%$. La diferencia es de
  una décima de punto porcentual.
  \\[0.3em]
  \textbf{Por qué importa}: el modelo normal reprodujo el dato real sin
  contar ni una ficha. Eso es lo que lo hace útil — permite evaluar cualquier
  corte, incluso uno que nunca se probó, usando sólo $\mu$ y $\sigma$.}

\ejercicio{3}{%
  Con los mismos $\mu=6{,}32$ y $\sigma=4{,}64$, calcula la puntuación $z$ si
  el servicio moviera el corte a 8, y otra vez si lo moviera a 13.
  ¿En cuál de los tres cortes (8, 10, 13) se marcaría a más gente?}{}
\renglones[4]
\respuesta{Corte 8: $z\approx 0{,}36$. Corte 13: $z\approx 1{,}44$.
Con 8 se marca a más gente.}
\solucion{%
  \[ z_{8} = \frac{8-6{,}32}{4{,}64} \approx 0{,}36 \qquad
     z_{13} = \frac{13-6{,}32}{4{,}64} \approx 1{,}44 \]
  \\[0.3em]
  Cuanto más bajo el corte, menor el $z$ y mayor el área a la derecha: se
  marca a más gente. Con corte 8 se marcaría aproximadamente al 36\% del
  archivo; con 13, alrededor del 7,5\%.
  \\[0.3em]
  Fíjate que no hizo falta volver a revisar ninguna ficha.}

\ejercicio{4}{%
  El servicio evalúa mover el corte. Con el corte en 10 tiene 88\% de
  sensibilidad, 88\% de especificidad y un VPP del 51,2\%.
  \begin{enumerate}[label=(\alph*),topsep=2pt,itemsep=1pt]
    \item Si baja el corte a 8, ¿qué le pasa a la sensibilidad? ¿Y a la
          especificidad? ¿Y al VPP?
    \item Si lo sube a 13, ¿qué cambia?
    \item El servicio tiene capacidad para 30 entrevistas al mes y tamiza 200
          estudiantes. ¿Qué corte recomendarías, y qué estás sacrificando?
  \end{enumerate}}{}
\renglones[8]
\respuesta{(a) Sensibilidad sube, especificidad baja, VPP baja.
(b) Al revés. (c) Respuesta abierta y argumentada.}
\solucion{%
  (a) Bajando el corte a 8 se marca a más gente: se escapan menos casos
  (\textbf{sensibilidad sube}) pero hay más falsas alarmas
  (\textbf{especificidad baja}), y como los nuevos marcados son en su mayoría
  sanos, el \textbf{VPP baja}.
  \\[0.3em]
  (b) Subiéndolo a 13 pasa exactamente lo contrario: sensibilidad baja
  (se escapan más casos), especificidad sube, VPP sube.
  \\[0.3em]
  \textbf{No existe un corte que mejore todo a la vez.} Todo corte es un
  intercambio entre los dos errores.
  \\[0.3em]
  (c) Con corte 10 se marcan 43 de 200 y hay capacidad para 30: el servicio ya
  está desbordado. Subir a 13 dejaría unos 15 marcados, que entran cómodos —
  pero se perderían casos reales, y el falso negativo es el error más grave en
  salud mental.
  \\[0.3em]
  Una respuesta defendible: mantener el corte en 10 y \textbf{priorizar} las
  entrevistas por puntaje (primero los más altos), en vez de subir el corte.
  Así no se descarta a nadie de antemano y se ordena la lista de espera por
  gravedad. Otra: pedir más recursos con este cálculo como argumento.
  \\[0.3em]
  \textit{Lo que se evalúa acá no es el número, es que puedan nombrar qué
  sacrifican con cada opción.}}
